\chapter{Specification and topology}
\label{ch:spec}

\section{Requirements}

\begin{table}[htbp]
\centering
\caption{Converter specification.}
\label{tab:spec}
\begin{tabular}{llr}
\toprule
Quantity & Symbol & Value \\
\midrule
DC-link voltage, nominal & $\Vdc$        & \SI{800}{\volt} \\
Rated phase current      & $I_\mathrm{rated}$ & \SI{30}{\Arms} \\
Switching frequency      & $\fsw$        & \SI{20}{\kilo\hertz} \\
Output fundamental       & $f_\mathrm{ref}$ & \SI{200}{\hertz} \\
Coolant temperature      & $T_\mathrm{liq}$ & \SI{65}{\celsius} (assumed) \\
Maximum ambient temperature & $T_\mathrm{amb}$ & \SI{50}{\celsius} \\
Maximum junction temperature & $\Tj^{\max}$ & \SI{150}{\celsius} \\
DC-link voltage ripple   & $\Delta V/\Vdc$ & $\leq \SI{0.5}{\percent}$ \\
\bottomrule
\end{tabular}
\end{table}

The output fundamental of \SI{200}{\hertz} corresponds to a four-pole machine
at \SI{6000}{rpm}, which is a representative traction operating point. The
coolant temperature of \SI{65}{\celsius} is \emph{not} given by the
specification: it is the value normally available from a vehicle's
power-electronics loop, a separate and hotter circuit than the one serving
the cabin, and it is adopted here as an assumption. Every thermal result is
therefore also given at \SI{70}{\celsius}.

\subsection{Requirements this design adds}

The specification asks for the rated current ``at any load power factor''.
It does not say at zero output frequency. This design adds that case
deliberately: a traction inverter holding a vehicle on a slope delivers
rated current into a stalled machine, and one phase then carries the full
\SI{42.4}{\ampere} peak as direct current. The rated current is
guaranteed at locked rotor, at full switching frequency and without
derating. It is a design choice stricter than the specification, and it is
the one that decides the device (\cref{sec:deviceselection}).

\subsection{Voltage hierarchy}
\label{sec:vhierarchy}

\begin{table}[htbp]
\centering
\caption{DC-link voltage levels and what each one sizes.}
\label{tab:vhierarchy}
\small
\begin{tabularx}{\textwidth}{lrY}
\toprule
Level & Value & Sizes \\
\midrule
Nominal, $\Vdc$                         & \SI{800}{\volt} & losses, operating point \\
Maximum continuous, $V_\mathrm{max,cont}$ & \SI{900}{\volt} & working voltage: creepage, static derating, DC-link \\
Overvoltage trip, $V_\mathrm{OV,trip}$   & \SI{950}{\volt} & comparator threshold; switching overshoot \\
Abnormal, $V_\mathrm{abn}$              & \SI{1000}{\volt} & withstand with the bridge disabled \\
\bottomrule
\end{tabularx}
\end{table}

Two derating rules are applied to the \SI{1200}{\volt} devices. In steady
state the blocking voltage stays below \SI{80}{\percent} of the rating at
$V_\mathrm{max,cont}$: $900/1200 = \SI{75}{\percent}$. During switching the
peak stays below \SI{90}{\percent}, \SI{1080}{\volt}, at $V_\mathrm{OV,trip}$,
which leaves \SI{130}{\volt} of overshoot budget. That peak is taken at
the die, not at the pins: the package inductance between the two adds
\SIrange{40}{75}{\volt} at turn-off, so a criterion verified at the pins
would pass while the die is already beyond it (\cref{sec:kelvin}).

$V_\mathrm{max,cont}$ is an assumption on the battery, not a property of
this module, and it is recorded in \cref{tab:sysassume}.

\subsection{Assumptions on the surrounding system}

The module is sized under four assumptions about what surrounds it. Each is
stated here once and used in the chapter that needs it.

\begin{table}[htbp]
\centering
\caption{Assumptions on the system around the power stage.}
\label{tab:sysassume}
\small
\begin{tabularx}{\textwidth}{YYY}
\toprule
Assumption & Value & Used in \\
\midrule
Maximum continuous bus voltage, including charging and regenerative braking & \SI{900}{\volt} & \cref{sec:vhierarchy}, \cref{sec:dclink}, \cref{sec:insulation} \\
Source impedance at \SI{20}{\kilo\hertz} & $L \geq \SI{1.5}{\micro\henry}$ or $R \geq \SI{50}{\milli\ohm}$ & \cref{sec:dclink}, \cref{ch:model} \\
Coolant at the rail inlet & 50/50 water--glycol, $\leq\SI{65}{\celsius}$, $\geq\SI{10}{\litre\per\minute}$, \SI{4.7}{\kilo\pascal} drop & \cref{sec:cooling} \\
Modulation at standstill & discontinuous (DPWM), recommended and not relied upon & \cref{ch:modulation}, \cref{sec:lrdpwm} \\
\bottomrule
\end{tabularx}
\end{table}

\section{Topology}

\subsection{Why a two-level voltage-source inverter}

\begin{table}[htbp]
\centering
\caption{Topology alternatives considered for an \SI{800}{\volt} traction drive.}
\label{tab:topology}
\small
\begin{tabularx}{\textwidth}{lYYY}
\toprule
 & \textbf{Two-level VSI} & Three-level NPC & Three-level T-type \\
\midrule
Devices                & \textbf{6} & 12 + 6 diodes & 12 \\
Device voltage class   & \SI{1200}{\volt} & \SI{650}{\volt} & \SI{650}{\volt} / \SI{1200}{\volt} \\
Output voltage levels  & 3 (line: 5) & 5 & 5 \\
Output harmonic content & highest & lower & lower \\
Gate drive channels    & \textbf{6} & 12 & 12 \\
Commutation loop       & \textbf{one per leg} & multiple, asymmetric & multiple \\
Neutral-point balancing & \textbf{not required} & required & required \\
\bottomrule
\end{tabularx}
\end{table}

\begin{figure}[htbp]
\centering
% Ponte trifase a due livelli: bus, DC-link, tre gambe con diodo di corpo, carico RL a stella.
% Riferimenti dei MOSFET come sulla scheda (Q1k1 lato alto, Q1k2 lato basso).
\begin{circuitikz}[font=\footnotesize, scale=0.9, transform shape]
  \def\yp{4.4}\def\ym{2.2}
  % bus
  \draw (0.3,\yp) node[left]{DC$+$} -- (11.2,\yp);
  \draw (0.3,0) node[left]{DC$-$} -- (11.2,0);
  % DC-link: film e ceramici
  \draw (1.3,\yp) to[C] (1.3,0);
  \draw (2.7,\yp) to[C] (2.7,0);
  \node[font=\scriptsize, align=right, anchor=east] at (1.15,3.3) {film\\\SI{200}{\micro\farad}};
  \node[font=\scriptsize, align=right, anchor=east] at (2.6,1.1) {C0G\\\SI{564}{\nano\farad}};
  % tre gambe
  \foreach \x/\k/\f in {4.6/1/a, 7.2/2/b, 9.8/3/c} {
    \draw (\x,\yp) -- (\x,3.95) node[nmos, bodydiode, anchor=D] (H\k) {};
    \draw (H\k.S) -- (\x,\ym);
    \draw (\x,0.45) node[nmos, bodydiode, anchor=S] (L\k) {};
    \draw (L\k.S) -- (\x,0);
    \draw (L\k.D) -- (\x,\ym);
    \node[anchor=east, font=\scriptsize] at ($(H\k.G)+(-0.1,0.25)$) {Q1\k1};
    \node[anchor=east, font=\scriptsize] at ($(L\k.G)+(-0.1,0.25)$) {Q1\k2};
    % uscita di fase: scende attraversando il bus negativo senza collegarsi
    \draw (\x,\ym) to[short, *-] (\x+0.7,\ym) -- (\x+0.7,0.12);
    \draw[white, line width=3pt] (\x+0.7,0.1) -- (\x+0.7,-0.1);
    \draw (\x+0.7,0.12) -- (\x+0.7,-0.6) node[right]{$\f$};
    \draw (\x+0.7,-0.6) to[R, l_=$R$] (\x+0.7,-2.2) to[cute inductor, l_=$L$] (\x+0.7,-3.6);
  }
  \draw (5.3,-3.6) -- (10.5,-3.6);
  \node[font=\scriptsize, align=left, anchor=west] at (10.8,-2.6) {motor, or the\\$RL$ bench load\\of \cref{ch:analytical}};
\end{circuitikz}

\caption{The two-level three-phase bridge with its DC-link: film capacitors for the
bulk energy, C0G ceramics next to each leg for the commutation loop. The
MOSFET references are those of the board; the reverse current flows through
the body diodes, since no external diode is fitted.}
\label{fig:vsi}
\end{figure}

The two-level bridge is chosen (\cref{fig:vsi}). Multilevel topologies exist
to reduce the voltage step seen by the output filter and to allow
lower-voltage, faster devices at high bus voltages. With a \SI{1200}{\volt}
SiC MOSFET the second reason is void: the device is fast enough to switch
\SI{800}{\volt} at \SI{20}{\kilo\hertz} with acceptable losses
(\cref{ch:losses}). The first remains real, but it costs twice the devices
and gate drives and a neutral-point balancing problem in the control
software. The two-level bridge also has \emph{one commutation loop per
leg}, geometrically simple and short; \cref{ch:layout} shows that loop
inductance is the binding constraint of the layout.

\subsection{Why silicon carbide}

\begin{table}[htbp]
\centering
\caption{Device technology comparison at \SI{800}{\volt}, \SI{20}{\kilo\hertz}.}
\label{tab:technology}
\small
\begin{tabularx}{\textwidth}{lYY}
\toprule
 & \textbf{SiC MOSFET} & Si IGBT (\SI{1200}{\volt}) \\
\midrule
Conduction characteristic & resistive, $\Rdson$ & voltage offset $V_\mathrm{CE(sat)}$ \\
Conduction loss at low current & \textbf{low} & offset dominates \\
Switching energy & \textbf{low} & tail current \\
Practical $\fsw$ & \SI{20}{}--\SI{50}{\kilo\hertz} & \SI{5}{}--\SI{15}{\kilo\hertz} \\
Reverse conduction & body diode, high $V_\mathrm{SD}$ & external diode required \\
Gate drive & \textbf{demanding}: negative bias, tight margins & tolerant \\
Cost & higher & lower \\
\bottomrule
\end{tabularx}
\end{table}

At \SI{20}{\kilo\hertz} the IGBT is excluded by its switching energy: the
tail current at turn-off makes the loss grow with frequency far faster than
in a MOSFET. With the fastest \SI{1200}{\volt},
\SI{40}{\ampere} TRENCHSTOP device~\cite{infineon_ikw40n120h3}, the
closed-form PWM loss formulas give \SI{643}{\watt} for the inverter at the
design point, \SI{69}{\watt} of switching loss per device, and a junction
temperature of \SI{151}{\celsius} before any overload. A hybrid with SiC
Schottky diodes removes the recovery contribution to $\Eon$; taking the
\emph{limit} $\Eon = 0$ still leaves \SI{332}{\watt}, and at locked rotor,
where the hottest device switches \SI{42.4}{\ampere} every period, it
reaches \SI{186}{\celsius}, beyond the absolute maximum. The margin
(factor 3--5 on the losses) is much wider than the uncertainty of an
analytical estimate, which is why no simulation was spent on it. The SiC MOSFET's resistive conduction characteristic is also
better matched to a traction duty cycle, where the converter spends most of
its life well below rated current, where an IGBT still pays its
saturation-voltage offset in full.

The price is paid at the gate: a low threshold voltage with a negative
temperature coefficient and edges measured in nanoseconds make the gate loop
and the commutation loop critical. The gate resistances and the dead time of
\cref{ch:losses} and the layout of \cref{ch:layout} follow from it.

\section{Device selection}
\label{sec:deviceselection}

The pass/fail criteria were fixed \emph{before} any result was computed, and
every candidate went through the same model, the same load and the same
thermal loop.

\subsection{Criteria, fixed before the results}

\begin{table}[htbp]
\centering
\caption{Pass/fail criteria for the power device.}
\label{tab:criteria}
\small
\begin{tabularx}{\textwidth}{lYY}
\toprule
 & Criterion & Why this threshold \\
\midrule
C1 & $\Tj \leq \SI{120}{\celsius}$ at locked rotor, hottest device, coolant
     \SI{65}{\celsius}, $\Rthcs = \SI{0.5}{\celsius\per\watt}$
   & the \SI{150}{\celsius} design limit less \SI{30}{\celsius} for the
     tolerances: maximum $\Rdson$ and $\Rthjc$, coolant at
     \SI{70}{\celsius}, pad at $+\SI{20}{\percent}$. Estimated here at
     about \SI{25}{\celsius}; verified all together, with the thermal loop
     closed, at \SI{29.4}{\celsius} (\cref{sec:worstcase}) \\
C2 & Peak $V_\mathrm{DS}$ \emph{at the die} $\leq \SI{1080}{\volt}$ at
     \SI{950}{\volt}, \SI{60}{\ampere}, \SI{20}{\nano\henry} external loop,
     at turn-off and on the opposite device during diode recovery
   & \cref{sec:vhierarchy}; \SI{60}{\ampere} is the overload trip, the
     largest current the device is ever asked to switch off in operation;
     \SI{20}{\nano\henry} is the layout budget, not the hoped-for value \\
C3 & Drain--source creepage on the board at \SI{900}{\volt}
   & PD1, \SI{2.8}{\milli\metre}~\cite{standardclarity} \\
C4 & AEC-Q101 (discretes) or AQG~324 (modules) & automotive application \\
C5 & Declared short-circuit withstand $t_\mathrm{SC} \geq \SI{1.5}{\micro\second}$
   & twice the typical desaturation chain of the gate driver
     (\SI{0.75}{\micro\second}); \num{1.7} times the worst-case one \\
\bottomrule
\end{tabularx}
\end{table}

\subsection{Candidates and results}

Candidates were taken in the \SI{1200}{\volt} TO-247-4L class from two
manufacturers (Wolfspeed~\cite{wolfspeed_e3m0040,wolfspeed_e3m0021,wolfspeed_e3m0016},
Infineon~\cite{infineon_aimzh040,infineon_aimzh020}), plus a paralleled pair,
a top-side-cooled package~\cite{wolfspeed_e3m0032u2} and two
modules~\cite{wolfspeed_ccb021,wolfspeed_edb010}. The Wolfspeed thermal
descriptions are the vendor's, converted to the format of the model and
checked value by value against the vendor's own converted file; the Infineon
ones were built from the datasheet graphs, since the vendor distributes them
only through a registered tool. C2 was simulated with the vendor SPICE
models, which exist for Wolfspeed and not, publicly, for Infineon.

\begin{table}[htbp]
\centering
\caption{Selection results. $R_g$ is given as turn-on/turn-off. Locked-rotor
(LR) temperatures at the gate resistance that satisfies C2 unless marked;
all from the same model, the same double-pulse bench and the same ideal cold
plate, which makes them comparable with one another; the design on the real
cooling rail is in \cref{sec:coolalt}.}
\label{tab:selection}
\small
\begin{tabular}{lcrrrl}
\toprule
Candidate & $R_g$ for C2 & Loss & $\Tj$ LR, \SI{65}{\celsius} & $\Tj$ LR, \SI{70}{\celsius} & Verdict \\
\midrule
40\,m$\Omega$, all vendors & --- & \SIrange{163}{178}{\watt} & \SIrange{140}{142}{\celsius}$^a$ & --- & fails C1 \\
E3M0032120K & --- & \SI{144}{\watt} & \SI{124.3}{\celsius}$^a$ & --- & fails C1 \\
E3M0021120K & \SI{7.5}{}/\SI{15}{\ohm} & $>$\SI{140}{\watt} & $>$\SI{120.7}{\celsius}$^c$ & $>$\SI{126.7}{\celsius}$^c$ & fails C1 \\
2 $\times$ E3M0040120K & $>\SI{25}{\ohm}$ & $>$\SI{150}{\watt} & $>$\SI{97.6}{\celsius}$^d$ & $>$\SI{103.2}{\celsius}$^d$ & fails C2 \\
CCB021M12FM3 & --- & \SI{154}{\watt} & \SI{153.1}{\celsius} & --- & fails C1, C4 \\
AIMZH120R020M1T & open & \SIrange{97}{128}{\watt} & \SIrange{97}{110}{\celsius}$^b$ & \SIrange{103}{116}{\celsius}$^b$ & C2 open \\
E3M0032120U2 (top side) & --- & --- & $>$\SI{124.3}{\celsius}$^e$ & --- & fails C1 \\
\textbf{\device{}} & \textbf{\SI{5}{}/\SI{12.5}{\ohm}} & \textbf{\SI{128}{\watt}} & \textbf{\SI{111.1}{\celsius}} & \textbf{\SI{116.8}{\celsius}} & \textbf{passes} \\
\bottomrule
\multicolumn{6}{l}{\footnotesize $^a$ at the datasheet $R_g$, before paying for C2; C1 already fails.
$^b$ over the window \SI{5}{}/\SI{7.5}{} to \SI{15}{}/\SI{20}{\ohm}.} \\
\multicolumn{6}{l}{\footnotesize $^c$ computed at \SI{5}{}/\SI{7.5}{\ohm}, already failing; the
\SI{7.5}{}/\SI{15}{\ohm} that C2 requires only adds loss.} \\
\multicolumn{6}{l}{\footnotesize $^d$ computed at \SI{15}{}/\SI{15}{\ohm}; excluded for the reasons in the
text, not for the overshoot figure.} \\
\multicolumn{6}{l}{\footnotesize $^e$ the same die in TO-247 already fails; the top pad is a quarter of
the tab, so $\Rthcs$ rises about \num{3.8} times.}
\end{tabular}
\end{table}

\begin{finding}{The first choice was neither automotive nor able to hold a stalled motor}
The C3M0040120K used in the first iterations of this work belongs to
Wolfspeed's industrial series; its automotive counterpart is the E3M
series. At locked rotor it reaches \SI{140.8}{\celsius}, twenty degrees
beyond C1, and the whole \SI{40}{\milli\ohm} class does the same whatever
the manufacturer. At the nominal point it had \SI{55}{\celsius} of margin,
which is why nobody had looked. Nor would discontinuous modulation have
saved it: with conduction dominating, clamping makes the hottest device
worse, and the best possible zero sequence still leaves it at
\SI{125}{\celsius} (\cref{sec:lrdpwm}).
\end{finding}

\begin{finding}{Overvoltage, not temperature, decided the selection}
With a TO-247 and \SI{20}{\nano\henry} of external loop, at
\SI{950}{\volt} and \SI{60}{\ampere}, neither of the two low-resistance
candidates keeps its die below \SI{1080}{\volt} at the datasheet
$R_g = \SI{2.5}{\ohm}$: the \device{} reaches \SI{1195}{\volt} at turn-off,
the E3M0021120K \SI{1208}{\volt}. The gate resistance that C2 demands costs
switching energy, and that energy re-opened the thermal question. The
E3M0021120K passes C1 at \SI{2.5}{\ohm} (\SI{111.6}{\celsius}) and already
fails it at \SI{5}{}/\SI{7.5}{\ohm} (\SI{120.7}{\celsius}), while C2 asks
for \SI{7.5}{}/\SI{15}{\ohm}. The \device{} needs \SI{5}{}/\SI{12.5}{\ohm}
and still passes C1 with \SI{8.9}{\celsius} to spare.

The turn-off resistance is set by the turn-off overshoot, the turn-on
resistance by the recovery overshoot on the opposite device; they are
separate constraints, and a split gate resistor meets each at the lowest
cost. Using \SI{12.5}{\ohm} at turn-on too would add
\SI{635}{\micro\joule} per cycle for no benefit
(\cref{tab:rgsweep}).
\end{finding}

Three other routes were closed. The paralleled pair passes C1 and is
rejected for three reasons, in order of weight. \emph{First}, it doubles the power stage: twelve devices, twelve gate
networks, and a gate resistor per die, which paralleled dies need to damp
the oscillation between their two gate loops. \emph{Second}, static
current sharing is self-correcting, because $\Rdson$ rises with
temperature, but dynamic sharing is not: the threshold of the
E3M0040120K spans \SIrange{1.8}{3.6}{\volt}~\cite{wolfspeed_e3m0040}, and
in a pair drawn at random the lower-threshold die turns on first and off
last, taking most of the switching energy of both. The simulated pair
shares within \SI{1}{\percent} only because its two model dies are
identical. \emph{Third}, one desaturation input has to watch two dies, and
two body diodes recover together through a longer and asymmetric loop,
which is why the bench needed \SI{25}{\ohm} to approach the limit at all.
The overshoot figure itself (\SI{1081}{\volt} against \SI{1080}{\volt}
at \SI{25}{\ohm}) is consistent with this, but it plays no part in the
rejection: one volt, on a model that misses the turn-on energy by a factor of 3.4
(\cref{sec:kelvin}), decides nothing.

Wolfspeed's top-side-cooled U2 package exists in the automotive range only
at \SI{32}{} and \SI{40}{\milli\ohm}, classes that already fail C1 in
TO-247, and its top pad is a quarter of the TO-247 tab, so the same
interface material gives about four times the case-to-sink resistance.
Top-side-cooled packages from other makers reach lower resistances; the
best of them is compared with the adopted solution on the real cooling rail
in \cref{sec:coolalt}, where it runs hotter. The modules
are either industrial and too hot (the six-pack) or meant to be sintered
onto a cooler (the automotive single switch), a manufacturing route this
design does not have.

\subsection{The choice}

Two candidates pass everything that can be demonstrated: the \device{} and
the AIMZH120R020M1T. Each leaves one criterion undemonstrated, a different
one in each case. Wolfspeed does not declare a
short-circuit withstand time (C5). Infineon declares it,
\SI{1.5}{\micro\second}, but publishes no SPICE model, so C2 cannot be
simulated; its thermal result holds across the whole range of gate
resistance a bench test could plausibly require.

The choice rests on one principle: a design is not adopted if the
verification that sizes it cannot be run. For this converter the sizing verification is C2: it
sets both gate resistors, and through them the switching loss, the
locked-rotor temperature and the dead time. With no public SPICE model for
the Infineon device, every one of those numbers would be an assumption
until a bench test existed. Better datasheet figures do not
change this; they are the reason the alternative is kept documented.

The \device{} is selected~\cite{wolfspeed_e3m0016}, because every
criterion that can be computed is computed with the manufacturer's own
models (thermal description and SPICE model alike), and the gate drive,
the footprint and the protection concept carry over unchanged.

One measurement would reverse the choice: a double pulse on the
AIMZH120R020M1T in the real power loop that keeps the die below
\SI{1080}{\volt} at a gate resistance inside the window of
\cref{tab:selection}. The Infineon device remains the documented
alternative: lower losses and a declared $t_\mathrm{SC}$, at the price of a
gate supply at \SI{+20}{\volt}, loss tables built by hand and an overshoot
figure that exists only after a bench test. The missing $t_\mathrm{SC}$ of
the chosen device is recorded as a residual risk in
\cref{app:assumptions}.

\begin{table}[htbp]
\centering
\caption{Main parameters of the \device{}~\cite{wolfspeed_e3m0016}. Switching
energies at \SI{800}{\volt}, \SI{80.28}{\ampere}, $R_{G(\mathrm{ext})} =
\SI{2.5}{\ohm}$, $\Tj = \SI{175}{\celsius}$.}
\label{tab:device}
\begin{tabular}{llrl}
\toprule
Parameter & Symbol & Value & Condition \\
\midrule
Breakdown voltage      & $V_\mathrm{(BR)DSS}$ & \SI{1200}{\volt} & \\
On-resistance          & $\Rdson$   & \SI{16}{\milli\ohm}  & \SI{25}{\celsius}, typ.; \SI{22}{} max. \\
                       &            & \SI{29}{\milli\ohm}  & \SI{175}{\celsius} \\
Total gate charge      & $Q_g$      & \SI{223}{\nano\coulomb} & \\
Input capacitance      & $C_\mathrm{iss}$ & \SI{6922}{\pico\farad} & \\
Reverse transfer cap.  & $C_\mathrm{rss}$ & \SI{13}{\pico\farad} & \\
Internal gate resistance & $R_{G(\mathrm{int})}$ & \SI{2.6}{\ohm} & \\
Gate threshold         & $V_\mathrm{GS(th)}$ & \SI{2.5}{\volt} & typ., \SI{25}{\celsius}; \SI{1.8}{} min. \\
Turn-on energy, SiC diode  & $\Eon$  & \SI{1287}{\micro\joule} & \\
Turn-off energy, SiC diode & $\Eoff$ & \SI{805}{\micro\joule} & \\
\textbf{Turn-on energy, body diode}  & $\Eon$  & \textbf{\SI{2552}{\micro\joule}} & \\
\textbf{Turn-off energy, body diode} & $\Eoff$ & \textbf{\SI{788}{\micro\joule}} & \\
Body-diode recovery charge & $Q_\mathrm{rr}$ & \SI{1665}{\nano\coulomb} & \SI{5.18}{\ampere\per\nano\second} \\
Junction-to-case resistance & $\Rthjc$ & \SI{0.23}{\celsius\per\watt} & typ.; \SI{0.31}{} max. \\
Pulsed drain current   & $I_\mathrm{D,pulse}$ & \SI{321}{\ampere} & \\
\bottomrule
\end{tabular}
\end{table}

Two entries in \cref{tab:device} govern the whole of \cref{ch:losses}.

First, the datasheet quotes switching energies for \emph{two different
free-wheeling arrangements}. In a bridge where the opposite device carries the
free-wheeling current, the relevant case depends on whether that device has an
antiparallel Schottky diode. This design has none, so reverse conduction goes
through the MOSFET body diode and the applicable turn-on energy is
\SI{2552}{\micro\joule}, not \SI{1287}{\micro\joule}: a factor of 2.0. The
difference is the reverse-recovery charge of the body diode.

Second, the package is the four-lead TO-247-4L, not the conventional
three-lead TO-247. The fourth lead is a Kelvin source connection that returns
the gate-drive current directly to the source of the die, bypassing the
package source inductance shared with the power path. \Cref{sec:kelvin}
quantifies what that pin is worth.
